% Options for packages loaded elsewhere
\PassOptionsToPackage{unicode}{hyperref}
\PassOptionsToPackage{hyphens}{url}
%
\documentclass[
  fleqn,compress,utf8,aspectratio=169,t,
]{beamer}

\usetheme[]{lmu}

\usepackage{pgfpages}
% Prevent slide breaks in the middle of a paragraph
\widowpenalties 1 10000
\raggedbottom

\makeatletter
\@ifundefined{KOMAClassName}{% if non-KOMA class
  \IfFileExists{parskip.sty}{%
    \RequirePackage{parskip}
  }{% else
    \setlength{\parindent}{0pt}
    \setlength{\parskip}{6pt plus 2pt minus 1pt}}
}{% if KOMA class
  \KOMAoptions{parskip=half}}
\makeatother

\IfFileExists{xurl.sty}{\usepackage{xurl}}{} % add URL line breaks if available
\IfFileExists{bookmark.sty}{\usepackage{bookmark}}{\usepackage{hyperref}}
\newif\ifbibliography
\urlstyle{same} % disable monospaced font for URLs

\usepackage{listings}
\lstset{defaultdialect=[5.3]Lua}
\lstset{defaultdialect=[x86masm]Assembler}
\lstset{
  basicstyle=\ttfamily
}
\usepackage{longtable,booktabs}
\usepackage{caption}
% Make caption package work with longtable
\makeatletter
\def\fnum@table{\tablename~\thetable}
\makeatother
\usepackage{graphicx}
\makeatletter
\def\maxwidth{\ifdim\Gin@nat@width>\linewidth.7\linewidth\else\Gin@nat@width\fi}
\def\maxheight{\ifdim\Gin@nat@height>\textheight.7\textheight\else\Gin@nat@height\fi}
\makeatother
% Scale images if necessary, so that they will not overflow the page
% margins by default, and it is still possible to overwrite the defaults
% using explicit options in \includegraphics[width, height, ...]{}
\setkeys{Gin}{width=\maxwidth,height=\maxheight,keepaspectratio}
% Set default figure placement to htbp
\makeatletter
\def\fps@figure{htbp}
\makeatother
\usepackage[normalem]{ulem}
% Avoid problems with \sout in headers with hyperref
\pdfstringdefDisableCommands{\renewcommand{\sout}{}}
\setlength{\emergencystretch}{3em} % prevent overfull lines
%Tightlist Command
\providecommand{\tightlist}{%
  \setlength{\itemsep}{0pt}\setlength{\parskip}{0pt}}

% Hacks for columns / blocks
\newcommand{\columnsbegin}[1]{\begin{columns}[t]%
    \column{#1}
    }
    \newcommand{\columnsend}{\end{columns}}
\newcommand{\blockbegin}[1]{\begin{block}{#1}}
    \newcommand{\blockend}{\end{block}}

\newcommand{\passthrough}[1]{#1}


%\beamertemplatenavigationsymbolsempty

\title{Title}
\subtitle{Subtitle}
\author{%
    Max Mustermann \and
    Karl Wombat}
\date{April 04, 2018}
\institute{MNM-Team, LMU München}

\usepackage{biblatex}
\addbibresource{lmu-example.bib}

\begin{document}
\begin{frame}
  \titlepage
\end{frame}

\hypertarget{overview}{%
\section{Overview}\label{overview}}

\begin{frame}{Motivation}
\protect\hypertarget{motivation}{}
\begin{itemize}
\tightlist
\item
  The LMU provides students, professors and other staff with

  \begin{itemize}
  \tightlist
  \item
    Corporate Design Manual~\autocite{lmu2006}
  \item
    Web-Styleguide~\autocite{lmu2007}
  \item
    templates for MS Office\footnote<.->{Although we really dislike MS
      Products}
  \end{itemize}
\item
  The IT-Service for Physics provides some \LaTeX~templates such as
  beamer theme~\autocite{lmu2014}.
\item
  The MNM Team again modernized this and changed the ratio to 16:9.
\end{itemize}
\end{frame}

\begin{frame}{Usage}
\protect\hypertarget{usage}{}
\begin{itemize}
\item
  install the style-files in your project-directory

  \begin{itemize}
  \tightlist
  \item
    beamerthemeLMU.sty
  \item
    beamerinnerthemelmu.sty
  \item
    beamerouterthemelmu.sty
  \item
    beamercolorthemelmu.sty
  \item
    beamerfontthemelmu.sty
  \end{itemize}
\item
  Write simple Markdown documents like this
\item
  We will show you some nice features.
\end{itemize}
\end{frame}

\begin{frame}[fragile]{Code}
\protect\hypertarget{code}{}
\begin{lstlisting}[language=TeX]
\documentclass[
    compress,    % compresses navigation bar
]{beamer}

\usetheme{LMU}   % enable LMU theme
\end{lstlisting}
\end{frame}

\begin{frame}{Description Lists}
\protect\hypertarget{description-lists}{}
You may consult further advice online. Note that this is a description
list

\begin{description}
\tightlist
\item[pandoc]
\url{https://pandoc.org/MANUAL.html}
\item[beamer]
\url{https://pandoc.org/MANUAL.html\#producing-slide-shows-with-pandoc}
\end{description}
\end{frame}

\begin{frame}{Block Environment}
\protect\hypertarget{block-environment}{}
\begin{block}{This is a block}
\protect\hypertarget{this-is-a-block}{}
This is the content in a block. This is caused by the slide-level
parameter in the YAML config. See above.
\end{block}
\end{frame}

\begin{frame}{General advice}
\protect\hypertarget{general-advice}{}
\begin{itemize}
\tightlist
\item
  Stay simple in Markdown
\item
  Do not overuse this. Markdown and Beamer is certainly some magic.
\item
  Making simple and nice looking presentations has never been simpler.
\end{itemize}

BUT\ldots{}

\pause

\begin{itemize}
\tightlist
\item
  Do not misuse. Use Latex if you need more complex things.
\item
  Markdown is \emph{powerful} due to its simplicity.
\item
  It is done for \textbf{easy} text editing.
\end{itemize}
\end{frame}

\begin{frame}{Tables}
\protect\hypertarget{tables}{}
Latex tables are always a pain. In Markdown it is much easier.

\begin{longtable}[]{@{}ll@{}}
\toprule
Syntax & Description\tabularnewline
\midrule
\endhead
Header & Title\tabularnewline
Paragraph & Text\tabularnewline
\bottomrule
\end{longtable}
\end{frame}

\begin{frame}{Figures}
\protect\hypertarget{figures}{}
\begin{figure}
\centering
\includegraphics{../lmulogo.pdf}
\caption{This is the caption}
\end{figure}
\end{frame}

\begin{frame}[allowframebreaks]{References}
  \bibliographytrue
  \nocite{*}
  \printbibliography[heading=none]
\end{frame}


\end{document}
